\section{Related work}
\label{sec:related-work}

\paragraph{JEPA-style and self-supervised image representations.}
The Joint-Embedding Predictive Architecture family
\citep{lecun2022path, assran2023ijepa} predicts in representation space rather
than pixel space, motivated by the conjecture that nuisance-invariant
representations emerge from learning to predict the latents of related views.
DINO \citep{caron2021dino} and DINOv2 \citep{oquab2024dinov2} train
transformers via self-distillation to produce features that transfer
across image distributions and downstream tasks; CLIP
\citep{radford2021clip} and OpenCLIP \citep{cherti2023openclip} build
transferable representations via contrastive image--text alignment;
MAE \citep{he2022mae} reconstructs masked pixels and VICReg
\citep{bardes2022vicreg} regularises variance, invariance, and
covariance of the embedding. We read all of these as different routes
toward representations that are approximately invariant to many
photometric and framing nuisances; our probe asks whether that
invariance generalises to nuisance families held out from the probe's
training distribution. Our probe is
narrow: the frozen-backbone, single-step, stable-pair subset of the JEPA family
applied as an evaluation probe, not a pretraining objective.

\paragraph{Domain-specific medical foundation models.}
BiomedCLIP \citep{zhang2023biomedclip} is a CLIP-style model trained on
PMC-15M, $\sim$$15$\,M biomedical image--text pairs from $\sim$$4.4$\,M
scientific articles. DermLIP \citep{yan2025derm1m} is a CLIP-style model trained on
Derm1M, a $\sim$1 million dermatology image--text corpus that aggregates
publicly available dermoscopy datasets and dermatology-website content.
A broader dermatology foundation-model line includes MONET
\citep{kim2024monet} (image--text alignment over dermatology textbook
figures and PubMed) and PanDerm \citep{yan2024panderm} (a multi-modal
self-supervised model trained on $\sim$$2$M dermatology images);
RETFound \citep{zhou2023retfound} is a comparable medical-domain
foundation-model precedent in retinal imaging, cited here for the
domain-specific FM design pattern rather than as a methodological
template for our probe. We use BiomedCLIP and DermLIP as third-party frozen weights,
not as artefacts to evaluate in their own right; the paper makes no
claim about either model beyond the test $\auroc$ each produces under
our specific probe.

\paragraph{Distribution-shift evaluation.}
WILDS \citep{koh2021wilds} and the in-search-of-lost-DG framing
\citep{gulrajani2021dg} establish the modern protocol for distribution-shift
evaluation. Our contribution is complementary: a deterministic, leakage-
controlled, three-family-disjoint nuisance design that lets the experimenter
precisely control the operation-level overlap between train and test nuisance.
Where WILDS varies natural distributions (hospital, scanner, country), we vary
synthetic nuisance families partitioned along the augmentation axis. The two
designs answer different questions; the methodology proposed here is most
useful when the natural axis of variation is hard to disentangle from the
target signal, as in dermoscopic re-photography. ISIC challenge artefacts
\citep{codella2018isic, combalia2019bcn20000} provide additional dermoscopy
corpora that remain unprobed in this work and are candidates for the EXP-009
partition pretraining set (\S\ref{sec:followup}).
